\chapter{Conclusions}
\label{ch:conclusions}

\section{Summary}

The two readings lead to a precise correspondence between evolution operators. In the electromagnetic reading, replacing the constant
shift by $G\circ\curl$ recovers Maxwell evolution in the zero-current sector when Gauss's laws hold; the six-component momentum records the
spatial currents, and the completed eight-component field also records the charge densities (\cref{ch:curl}). A second construction, on a
constant-rate background, has a massive two-form sector with positive canonical energy (\cref{ch:mass}). Coupled to a connection it keeps
its constraint structure (\cref{ch:bcoupling}), but its tree amplitudes grow with energy, in transverse and longitudinal channels. This
motivates a Yang--Mills--Higgs completion, whose vector and scalar interactions cancel the leading growth of the tree amplitudes studied here
at fixed nonexceptional angles (\cref{ch:quantum,ch:connection}).

With the spectator Lorentz action, Clifford transport sends the temporal clock generator to the ambient volume element. At one constituent it
identifies the temporal Clifford algebra with the Lorentz commutant; inside that commutant the admissible symmetric involution for the third
element of the rotation triple is fixed up to sign, and the chosen orientation selects $\Gamma$. The transported volume $\sigma_1(T)=-\Omega$ acts by left multiplication on the embedded field-strength bivectors and
implements $G$. This relates the generators of the two rotations of Part~I; equating their physical frequencies further requires a common
time parameter and the normalisation $|y|v=\kappa$ (\cref{ch:hyp}). At general rung $n$ the Lorentz commutant is $\Cl(3n,3n-3)$ and the
transported rotations form $\so{3n}$ (\cref{ch:rungs}).

The second rung gives the principal matter result. With the specified temporal and spatial rotations gauged, the full $\R^{64}$ spinor module
carries one chiral Pati--Salam family with a right-handed neutrino, with cancelling perturbative anomalies and an even number of doublets for
each weak factor. For the chosen Standard-Model embedding the invariant mass kernels form one complex right-handed Majorana line, which does
not come from the transported temporal algebra. A nonzero value on this line, together with the bidoublet prescribed by the generator of
mass, leaves colour and $U(1)_Q$ unbroken; this identifies the electromagnetic generator and its unbroken abelian gauge direction in that phase.
Whether the gauge field of the first rung continues dynamically into that photon remains open, and selecting a stable vacuum with those
values, including the relative orientation of the two condensates, remains a dynamical question (\cref{ch:rungs,ch:ew}).

The restricted neutral Dirac ansatz gives $z_u-z_d-z_\nu+z_e=0$. For the allowed quartic $\Tr[(mm^\top)^2]$ the gradient component along the
omitted direction $D$ is $-\half v_R^2\kappa$ at $v_R\Phi_0+\kappa\,G_2\,\omega_S$. The adopted symmetries thus allow a term with a nonzero gradient
transverse to the restricted ansatz at this background, which generically sources a value along $D$; the size of the displacement depends on
the other terms and on the restoring curvature. An independent coefficient $d$ changes the alternating sum to $8d$, and the four real Dirac
coefficients can then parametrise any four signed entries of one family. Determining their values and the observed flavour structure
requires further dynamics (\cref{thm:potential,thm:splitgen}).

\section{Assumptions, obstructions and open problems}
\label{sec:open}

The following list collects, topic by topic, what the results assume, what was found to be obstructed, and what remains open.

\begin{enumerate}[leftmargin=2em,itemsep=4pt]
\item \emph{The hypothesis and the third rung.} Clifford transport, (A1) to (A4) of \cref{def:hyp}, is adopted and not derived. It includes
the gauging of the spatial directions $\mathfrak s_n$ of the other constituents and the spectator Lorentz action, in which the four-plane and
$\lint$ are those of the first constituent. At the first rung it gives the Pauli--G\"ursey algebra on a single Weyl doublet, which carries
the global anomaly as a quantum gauge theory (\cref{prop:onedoublet}). At the second rung it gives one family. At the third rung
$G_3=\sigma_3(T_3)$ is defined and is an involution, with no rate; the gauge algebra $\so9\oplus\so8$ acts on the irreducible
$(\mathbf{16},\mathbf8)$ and permits a common mass for every mode (\cref{thm:massrung,prop:rung3s}). Open are the interaction that the third
rung carries; whether a physical rest rotation there needs an alternative assignment that would supplement or replace (A3), such as the
image of $T_1$, a rotation that keeps only $\so3\oplus\so6$, or the central $-\Omega$, which is not the image of a temporal element; and the
behaviour of all results at rungs two and three under the mixed action of \PI{Section~5.8} instead of the spectator action.
\item \emph{The massive two-form and its completion.} The choice of the positive half excludes the duality rotations as symmetries of that
half with the adopted canonical energy (\cref{prop:mirror}); the free equations conserve other quantities, positive and duality invariant.
The massive equation has no first-order factorisation of the form considered on the octants (\cref{prop:nofirst}). Equal squared rates are
necessary for the copy transformations to be exact symmetries of the prescribed mass term (\cref{prop:equalrates}), and one clock is the
assumption that enforces them. A prescribed and a dynamical clock have different solution sets, and the elimination of $\bE$ with a dynamical
clock can be singular (\cref{thm:clock,prop:clockhalf}). The scattering results delimit the perturbative model: with a connection or a
scalar clock its tree amplitudes grow (\cref{thm:compton,thm:clockamp}), and the completion cancels the growth at fixed nonexceptional angles
with an assumed Yang--Mills term, kinetic term and potential of the Higgs vector (\cref{thm:bounded}); the ultraviolet theory is open.
\Cref{thm:ymforced} fixes the quadratic Lagrangian of the first-rung curvature up to its scale and a topological term, under the invariances
it assumes; at the second rung the direct product permits independent gauge couplings. Open are whether the split form of \cref{ch:curl},
with its additional modes, and the density of \cref{prop:BP}, which has both branches, have a role in the coupled theory; whether the clock
angle $\theta$ is a field of the completed theory, which would need a coupling that \cref{eq:Lcomp} does not contain; and everything beyond
tree level.
\item \emph{The vacuum.} The Pati--Salam direction $\Phi_0$ is selected by what it keeps (\cref{thm:psstage}) and the electroweak direction by
(A3). \Cref{prop:ewdir} identifies a critical direction on the fixed-norm sphere for potentials that preserve the gauge algebra and the global
phases; it does not select a radial stationary point or a stable global minimum. A quartic of the right triplet selects the orbit of
$\Phi_0$ within that component, while the trace quartic alone does not distinguish the right triplet from the left one; separate quartics do
not fix the relative orientation of the triplet and a rank-one bidoublet, which is physical and can break $Q$ (\cref{prop:select}). The
scales $v_R$ and $\kappa$ are free; the decoupling regime $v_R\gg|\kappa|$, in which the stated mass ratio approaches $\tfrac58$ and $\rho$
approaches one, is a tuning, and the coefficients $g$, $Z$, $\lambda$, $y$ and
the identifications $v=\kappa$ and $y=1$ of \cref{sec:coupled} are assumptions. The Majorana kernel has the form $\Phi_0=C\omega_S$ with $C$ in the full
Lorentz commutant; diagonal weak invariance excludes every such representation with $C$ in the transported temporal algebra
(\cref{prop:notransport}). Open are a distinguished construction of $\Phi_0$ from the additional spatial structure, or a potential that
selects the joint vacuum, stable in every direction, and fixes its relative scales; how the conversion $c$ of \PI{Section~9.6} acts on the
mass field; and what fixes the rate $\kappa$ (\PI{Prop.~4.49}).
\item \emph{The couplings.} The matching value $\sin^2\theta_W=\tfrac38$ holds if the three couplings of $\su4\oplus\su2_L\oplus\su2_R$ are
equal, at a scale that is not specified. The limiting vector-mass ratio $m_W^2/m_Z^2\to\tfrac58$ in the stated vacuum as $v_R/\kappa\to\infty$
is a separate statement (\cref{thm:vectors}). Their relation at low energy requires running and threshold information, which is not
computed.
\item \emph{The masses of a family.} The generator of mass alone gives one Dirac entry (\cref{thm:fermionmasses}); the Yang--Mills values split
it subject to a sum rule within the restricted ansatz (\cref{thm:split}); the sum rule is not protected (\cref{thm:potential}); and the general
split has four free values per family (\cref{thm:splitgen}). The observed hierarchy of the top quark over the bottom quark and the tau
motivates a near-pure-up interpretation of the heaviest family within the restricted ansatz. The neutrino Dirac entry is undetermined: with
an independent $d$ the same charged masses admit different neutrino Dirac entries, and the data do not locate the vacuum in the full neutral
space. Open are what sets $\kappa$, $x_3$, $x_B$ and $d$, and the size of $d$, which depends on the couplings and on the curvature of the
potential.
\item \emph{The number of families.} The second rung holds one family. With the interaction of the second rung held, the third rung gives
eight chiral copies; on restriction to the embedded Pati--Salam algebra built with $\Gamma_3$ it gives four families and four mirrors, which
the full third-rung algebra mixes (\cref{thm:families}). Four routes to three families are retained as alternatives with different
prerequisites, and none is chosen. (a) The third rung with the interaction of the second held: three families would be one triplet of the
Fock grading, the octants of weight one, where Part~I places $\bE$; this needs a selection of the light sector, compatible with the
interactions retained, by which only that weight carries light matter, since the invisible width of the $Z$ constrains additional
kinematically accessible light neutrinos with the standard coupling to the $Z$ \cite{PDGnu}. (b) Three constituents as three generations, the
reading of \PI{Section~10.1.6}: unequal constituent rates restrict an unbroken copy symmetry in the chosen model (\cref{prop:equalrates}), but
they do not decide whether a broader family symmetry could be gauged and spontaneously broken, and a construction that labels the families by
constituents is not provided by \cref{thm:families}. (c) Triality: the third rung gauges $\so8$ (\cref{prop:rung3s}), whose triality permutes
three inequivalent eight-dimensional representations of $\mathrm{Spin}(8)$. This is an analogy motivated by the conjecture of
\cite{Boyle2020}, which links triality and the exceptional Jordan algebra to three generations; a third-rung $\mathrm{Spin}(8)$ action alone
does not realise that construction, the automorphism supplies no three chiral families, and producing them in this tower requires an
explicit selection mechanism beyond the hypothesis. (d) The family of the second rung taken three times, as the Standard
Model does, which leaves the number three unexplained.
\item \emph{The mixings, the neutrinos and the uneaten direction.} With diagonal family values each charged fermion can be given its measured
mass (\cref{ex:observed}), and the family mass matrices are then aligned: mixing requires a misalignment of the relevant family mass matrices,
and whether the route to the families allows one that reproduces the measured mixings together with the masses is open. The independent
inputs of the neutrino illustration are the normal ordering, $m_1=0$, the two measured splittings and the chosen Majorana scale; the other two
light masses follow from them. If the full classical action preserves the proposed global phase, the
chosen electroweak vacuum has one broken phase direction independent of the gauge orbit, and its current has a colour anomaly; a
gauge-invariant term such as $\mu^2\Re\det F$, for the complex bidoublet $F$, can also break the phase explicitly and contribute to its mass
at tree level (\cref{thm:uneaten}). Its mass and couplings therefore remain undetermined.
\end{enumerate}
